% !TEX program = xelatex
% ============================================================================
%  POSN Computer Qualification — Number Theory
%  Topic: ทฤษฎีจำนวน (Number Theory)
% ============================================================================

\documentclass[11pt, aspectratio=169]{beamer}
% ============================================================================
%  POSN Computer Qualification — Shared Beamer Preamble
%  Used by all topic slide decks. Compile with XeLaTeX.
% ============================================================================

% ---------- Thai language & font ----------
\usepackage{iftex}
\usepackage{fontspec}

% XeTeX-specific: enable Thai line breaking
\ifXeTeX
  \XeTeXlinebreaklocale{th}
\fi
% Noto Sans Thai — body text (clean modern sans-serif, handles all Thai)
\setmainfont{Noto Sans Thai}[
  UprightFont = Noto Sans Thai Light,
  BoldFont    = Noto Sans Thai Medium,
  Scale       = 0.9
]
\setsansfont{Noto Sans Thai}[
  UprightFont = Noto Sans Thai Light,
  BoldFont    = Noto Sans Thai Medium,
  Scale       = 0.9
]

% Noto Sans Thai — clean modern sans-serif for structural elements
\newfontfamily\headingfont{Noto Sans Thai}[
  UprightFont = Noto Sans Thai Medium,
  BoldFont    = Noto Sans Thai Bold,
  Scale       = 0.9
]

% --- Heading font for structural elements ---
\setbeamerfont{title}{family=\headingfont, series=\bfseries, size=\huge}
\setbeamerfont{subtitle}{family=\headingfont, series=\mdseries}
\setbeamerfont{author}{family=\headingfont, series=\mdseries}

% --- Custom Title Page ---
\defbeamertemplate{title page}{custom}[1][]{%
  \centering
  \vspace*{2cm}
  % Title — in white
  {\usebeamerfont{title}\color{white}\inserttitle\par}
  \vspace{0.6cm}
  % Subtitle — in white
  {\usebeamerfont{subtitle}\color{white}\insertsubtitle\par}
  \vspace{2.5cm}
  % Author — blue filled rounded box, white text
  \begin{center}
    \begin{tcolorbox}[arc=3mm, colback=boxBlue, colframe=boxBlue!80,
                      width=0.42\textwidth, halign=center,
                      left=1.5em, right=1.5em, top=0.8em, bottom=0.8em,
                      boxrule=1.5pt, coltext=white]
      \usebeamerfont{author}\insertauthor
    \end{tcolorbox}
  \end{center}
}
\setbeamertemplate{title page}[custom]
\setbeamerfont{frametitle}{family=\headingfont, series=\bfseries}
\setbeamerfont{section in toc}{family=\headingfont}
\setbeamerfont{page number in head/foot}{family=\headingfont}
\setbeamerfont{section title}{family=\headingfont, series=\bfseries}
\setbeamerfont{subsection title}{family=\headingfont}

% Monospace font (for code/verbatim)
\setmonofont{Latin Modern Mono}[Scale=0.9]

% ---------- Math ----------
\usepackage{amsmath, amssymb}

% ---------- Graphics ----------
\usepackage{graphicx}
\usepackage{tikz}

% ---------- String parsing (for section headers) ----------
\usepackage{xstring}

% ---------- Colored boxes ----------
\usepackage[skins, breakable, xparse]{tcolorbox}
% Note: we do NOT load the 'most' option — it predefines example/solution environments
% that would conflict with our custom boxes.

% --- Color definitions ---
\definecolor{boxBlue}{RGB}{41, 128, 185}
\definecolor{boxGreen}{RGB}{39, 174, 96}
\definecolor{boxOrange}{RGB}{230, 126, 34}
\definecolor{boxPurple}{RGB}{142, 68, 173}
\definecolor{boxBlueBg}{RGB}{235, 245, 255}
\definecolor{boxGreenBg}{RGB}{235, 255, 240}
\definecolor{boxOrangeBg}{RGB}{255, 245, 235}
\definecolor{boxPurpleBg}{RGB}{248, 240, 255}

% --- Explanation box (blue) — for definitions, concepts, notes ---
\newtcolorbox{explanation}[1][]{
  enhanced,
  breakable,
  arc         = 2mm,
  colback     = boxBlueBg,
  colframe    = boxBlue!50,
  title       = #1,
  fonttitle   = \bfseries\color{white},
  coltitle    = white,
  colbacktitle= boxBlue,
  attach boxed title to top left = {yshift=-2mm, xshift=2mm},
  boxed title style = {arc=2mm, size=small},
  before skip = 8pt,
  after skip  = 8pt
}

% --- Example box (green) — for worked examples ---
\newtcolorbox{exambox}[1][]{
  enhanced,
  breakable,
  arc         = 2mm,
  colback     = boxGreenBg,
  colframe    = boxGreen!50,
  title       = #1,
  fonttitle   = \bfseries\color{white},
  coltitle    = white,
  colbacktitle= boxGreen,
  attach boxed title to top left = {yshift=-2mm, xshift=2mm},
  boxed title style = {arc=2mm, size=small},
  before skip = 8pt,
  after skip  = 8pt
}

% --- Solution box (orange) — for solutions and answers ---
\newtcolorbox{solbox}[1][]{
  enhanced,
  breakable,
  arc         = 2mm,
  colback     = boxOrangeBg,
  colframe    = boxOrange!50,
  title       = #1,
  fonttitle   = \bfseries\color{white},
  coltitle    = white,
  colbacktitle= boxOrange,
  attach boxed title to top left = {yshift=-2mm, xshift=2mm},
  boxed title style = {arc=2mm, size=small},
  before skip = 8pt,
  after skip  = 8pt
}

% --- Intuition box (purple) — for plain-language explanations, analogies ---
\newtcolorbox{intuition}[1][]{
  enhanced,
  breakable,
  arc         = 2mm,
  colback     = boxPurpleBg,
  colframe    = boxPurple!50,
  title       = #1,
  fonttitle   = \bfseries\color{white},
  coltitle    = white,
  colbacktitle= boxPurple,
  attach boxed title to top left = {yshift=-2mm, xshift=2mm},
  boxed title style = {arc=2mm, size=small},
  before skip = 8pt,
  after skip  = 8pt
}

% ---------- Beamer theme ----------
\usetheme{Madrid}
\usecolortheme{default}

% --- Custom beamer colors (blue accent matching explanation boxes) ---
\setbeamercolor{structure}{fg=boxBlue}
\setbeamercolor{palette primary}{fg=white, bg=boxBlue}
\setbeamercolor{palette secondary}{fg=white, bg=boxBlue!80}
\setbeamercolor{palette tertiary}{fg=white, bg=boxBlue!60}
\setbeamercolor{block title}{fg=white, bg=boxBlue}
\setbeamercolor{block body}{fg=black, bg=boxBlueBg}

% Remove navigation symbols for clean look
\setbeamertemplate{navigation symbols}{}

% Note: Frame title prefix removed — \addtobeamertemplate conflicts with beamer internals

% --- Outline (Table of Contents) styling ---
\setbeamerfont{section in toc}{family=\headingfont, series=\mdseries}
\setbeamerfont{subsection in toc}{family=\headingfont, series=\mdseries}
\setbeamertemplate{section in toc}{%
  \leavevmode\leftskip=1.5em
  \textcolor{boxBlue}{\textbf{\large\inserttocsectionnumber.}}%
  \quad\inserttocsection\par
  \vspace{0.2em}
}
\setbeamertemplate{subsection in toc}{%
  \leavevmode\leftskip=4em
  \textcolor{boxBlue!60}{\small$\bullet$}%
  \quad\inserttocsubsection\par
  \vspace{0.1em}
}

% Section header slides — Thai in white, English in smaller light blue below
\AtBeginSection{
  \begin{frame}[plain]{}
    \vspace*{2cm}
    \begin{beamercolorbox}[sep=12pt, center, rounded=true, shadow=true]{palette primary}
      \IfSubStr{\secname}{(}{%
        \StrBefore{\secname}{(}[\sectionthai]%
        \StrBetween[1,1]{\secname}{(}{)}[\sectioneng]%
        \begin{tabular}{@{}c@{}}
          {\usebeamerfont{title}\sectionthai} \\[0.2em]
          {\usebeamerfont{subtitle}\color{boxBlue!60}\sectioneng}
        \end{tabular}%
      }{%
        {\usebeamerfont{title}\secname\par}%
      }
    \end{beamercolorbox}
    \vfill
  \end{frame}
}

% Subsection header slides — smaller, centered
\AtBeginSubsection{
  \begin{frame}[plain]{}
    \vfill
    \centering
    {\usebeamerfont{section title}\insertsubsectionhead\par}
    \vfill
  \end{frame}
}

% Minimal footer: page number only, right-aligned
\setbeamertemplate{footline}{
  \hfill\usebeamerfont{page number in head/foot}\insertframenumber{} / \inserttotalframenumber\hspace*{1em}\vspace*{0.5ex}
}

% ---------- Spacing & readability ----------
\linespread{1.2}

% ---------- Useful commands ----------

% Highlight important text in blue
\newcommand{\highlight}[1]{\textcolor{boxBlue}{\textbf{#1}}}

% Math set notation helpers
\newcommand{\R}{\mathbb{R}}
\newcommand{\N}{\mathbb{N}}
\newcommand{\Z}{\mathbb{Z}}
\newcommand{\Q}{\mathbb{Q}}

% ============================================================================
%  End of preamble
% ============================================================================


\title[ทฤษฎีจำนวน]{\textcolor{white}{ทฤษฎีจำนวน} \textcolor{boxBlue!60}{(Number Theory)}}
\subtitle{การหารลงตัว • จำนวนเฉพาะ • ห.ร.ม./ค.ร.น. • สมภาค (Modular Arithmetic) }
\author[None W. (Yuan)]{None W. (Yuan)}
\date{}

\begin{document}

\begin{frame}[plain]\titlepage\end{frame}
\begin{frame}{สารบัญ (Outline)}\tableofcontents\end{frame}

\begin{frame}{ก่อนเรียน (Prerequisites)}
  \begin{explanation}[สิ่งที่ควรรู้ก่อนเรียน]
    \begin{itemize}
      \item \textbf{จำนวนเต็ม} — การบวก ลบ คูณ หาร, การหารเอาเศษ
      \item \textbf{เลขยกกำลัง} — $a^n$, การคูณเลขยกกำลัง
      \item \textbf{จำนวนจริง} — ความรู้พื้นฐานเกี่ยวกับจำนวน
    \end{itemize}
  \end{explanation}

  \begin{intuition}[ทฤษฎีจำนวนคืออะไร?]
    การศึกษาเกี่ยวกับ \textbf{จำนวนเต็ม} และสมบัติของมัน —
    โดยเฉพาะการหารลงตัว จำนวนเฉพาะ และการหาเศษจากการหาร
  \end{intuition}
\end{frame}

% ===========================================================================
%  SECTION 1: การหารลงตัว
% ===========================================================================
\section{การหารลงตัว (Divisibility)}

\begin{frame}{การหารลงตัว}
  \begin{explanation}[นิยาม]
    $a$ \textbf{หาร} $b$ \textbf{ลงตัว} เขียนแทนด้วย $a \mid b$
    ก็ต่อเมื่อ มีจำนวนเต็ม $k$ ที่ทำให้ $b = ak$

    \textbf{อ่านว่า}: ``$a$ divides $b$'' หรือ ``$a$ เป็นตัวหารของ $b$''
  \end{explanation}

  \begin{exambox}[ตัวอย่าง]
    $3 \mid 12$ เพราะ $12 = 3 \times 4$ \qquad $7 \mid 0$ เพราะ $0 = 7 \times 0$ \qquad $5 \nmid 13$
  \end{exambox}

  \begin{explanation}[ขั้นตอนการหาร (Division Algorithm)]
    สำหรับ $a \in \Z$, $b > 0$: มี $q, r$ เพียงคู่เดียวที่ \highlight{$a = bq + r$}, $0 \leq r < b$

    $q$ = ผลหาร, $r$ = เศษ

    \textbf{ตัวอย่าง:} $17 = 5 \times 3 + 2$ \qquad $-13 = 5 \times (-3) + 2$ \quad (เศษ $\geq 0$ เสมอ!)
  \end{explanation}
\end{frame}

\begin{frame}{กฎการหารลงตัว (Divisibility Rules)}
  \begin{columns}[T]
    \column{0.5\textwidth}
      \begin{explanation}[หลักหน่วย/หลักสุดท้าย]
        \textbf{2}: หลักหน่วย $0,2,4,6,8$ \\
        \textbf{5}: หลักหน่วย $0$ หรือ $5$

        \textbf{10}: หลักหน่วย $0$ \\
        \textbf{4}: 2 หลักสุดท้าย $\div 4$ ลงตัว

        \textbf{8}: 3 หลักสุดท้าย $\div 8$ ลงตัว
      \end{explanation}
    \column{0.5\textwidth}
      \begin{explanation}[ผลรวมเลขโดด]
        \textbf{3}: ผลรวมเลขโดด $\div 3$ ลงตัว

        \textbf{9}: ผลรวมเลขโดด $\div 9$ ลงตัว
      \end{explanation}
  \end{columns}

  \begin{exambox}[ตัวอย่าง]
    $12345$: หารด้วย $5$ ลงตัว (หลักหน่วย $5$), หารด้วย $3$ ลงตัว ($1+2+3+4+5=15$)

    $2024$: หารด้วย $4$ ลงตัว ($04 \div 4 = 1$)
  \end{exambox}
\end{frame}

% ===========================================================================
%  SECTION 2: จำนวนเฉพาะและจำนวนประกอบ
% ===========================================================================
\section{จำนวนเฉพาะและจำนวนประกอบ}

\begin{frame}{จำนวนเฉพาะ (Prime Numbers)}
  \begin{explanation}[นิยาม]
    \textbf{จำนวนเฉพาะ} คือจำนวนเต็มบวกที่มากกว่า $1$ และมีตัวหารบวกเพียงสองตัวคือ $1$ และตัวมันเอง

    \textbf{จำนวนประกอบ (Composite)} คือจำนวนเต็มบวกที่มากกว่า $1$ และไม่ใช่จำนวนเฉพาะ
  \end{explanation}

  \begin{columns}[T]
    \column{0.4\textwidth}
      \begin{exambox}[จำนวนเฉพาะ $< 50$]
        $2, 3, 5, 7, 11, 13,$

        $17, 19, 23, 29,$

        $31, 37, 41, 43, 47$

        \medskip
        \highlight{$2$ เป็นจำนวนเฉพาะคู่เพียงตัวเดียว!}
      \end{exambox}
    \column{0.6\textwidth}
      \begin{explanation}[หมายเหตุ]
        \begin{itemize}
          \item $1$ \textbf{ไม่ใช่} จำนวนเฉพาะ และ \textbf{ไม่ใช่} จำนวนประกอบ
          \item จำนวนเฉพาะมี \textbf{ไม่จำกัด}
        \end{itemize}
      \end{explanation}

      \begin{exambox}[ตัวอย่าง: จำนวนประกอบ]
        $4 = 2 \times 2$ \qquad $30 = 2 \times 3 \times 5$

        $91 = 7 \times 13$ \quad (91 ไม่ใช่จำนวนเฉพาะ!)
      \end{exambox}
  \end{columns}
\end{frame}

\begin{frame}{การแยกตัวประกอบเฉพาะ (Prime Factorization)}
  \begin{explanation}[ทฤษฎีบทหลักมูลของเลขคณิต]
    จำนวนเต็มบวกทุกตัวที่มากกว่า $1$ สามารถเขียนในรูปผลคูณของจำนวนเฉพาะได้
    \textbf{เพียงแบบเดียว} (ไม่นับลำดับ)
  \end{explanation}

  \begin{exambox}[ตัวอย่าง]
    $60 = 2^2 \times 3 \times 5$ \qquad $84 = 2^2 \times 3 \times 7$

    $126 = 2 \times 3^2 \times 7$ \qquad $1000 = 2^3 \times 5^3$
  \end{exambox}

  \begin{exambox}[การตรวจสอบว่าเป็นจำนวนเฉพาะ]
    ในการตรวจสอบว่า $n$ เป็นจำนวนเฉพาะหรือไม่
    ทดลองหารด้วยจำนวนเฉพาะทุกตัวที่ $\leq \sqrt{n}$ ก็พอ

    \textbf{ตัวอย่าง:} $101$ — ทดลอง $2, 3, 5, 7$ ($\sqrt{101} \approx 10$)

    ไม่มีตัวใดหารลงตัว $\rightarrow$ $101$ เป็นจำนวนเฉพาะ $\checkmark$
  \end{exambox}
\end{frame}

% ===========================================================================
%  SECTION 3: ห.ร.ม. และ ค.ร.น.
% ===========================================================================
\section{ห.ร.ม. และ ค.ร.น. (GCD \& LCM)}

\begin{frame}{ห.ร.ม. (GCD) และ ค.ร.น. (LCM)}
  \begin{explanation}[นิยาม]
    \textbf{ห.ร.ม.} ของ $a, b$ คือจำนวนเต็มบวกที่มากที่สุดที่หารทั้ง $a$ และ $b$ ลงตัว

    \textbf{ค.ร.น.} ของ $a, b$ คือจำนวนเต็มบวกที่น้อยที่สุดที่ $a$ และ $b$ หารลงตัว
  \end{explanation}

  \begin{columns}[T]
    \column{0.5\textwidth}
      \begin{exambox}[วิธีแยกตัวประกอบ]
        $12 = 2^2 \times 3$ \quad $18 = 2 \times 3^2$

        GCD $= 2^{\min(2,1)} \times 3^{\min(1,2)} = 2 \times 3 = 6$

        LCM $= 2^{\max(2,1)} \times 3^{\max(1,2)} = 2^2 \times 3^2 = 36$
      \end{exambox}
    \column{0.5\textwidth}
      \begin{explanation}[สูตรสำคัญ]
        \highlight{$ab = \text{GCD}(a,b) \times \text{LCM}(a,b)$}

        $12 \times 18 = 6 \times 36 = 216$ \quad $\checkmark$
      \end{explanation}
  \end{columns}
\end{frame}

\begin{frame}{ขั้นตอนวิธีของยุคลิด (Euclidean Algorithm)}
  \begin{explanation}[วิธีหา GCD โดยไม่ต้องแยกตัวประกอบ]
    \textbf{หลักการ}: $\text{GCD}(a,b) = \text{GCD}(b, a \bmod b)$ \quad (ทำซํ้าจนเหลือ $0$)
  \end{explanation}

  \begin{exambox}[ตัวอย่าง: Euclidean Algorithm (252 กับ 105)]
    $252 = 2 \times 105 + 42$ \quad $\rightarrow$ \quad GCD(105, 42)

    $105 = 2 \times 42 + 21$ \quad $\rightarrow$ \quad GCD(42, 21)

    $42 = 2 \times 21 + 0$ \quad $\rightarrow$ \quad \highlight{GCD = 21}
  \end{exambox}

  \begin{exambox}[ตัวอย่าง: Euclidean Algorithm (126 กับ 35)]
    $126 = 3 \times 35 + 21$ \quad $\rightarrow$ \quad GCD(35, 21)

    $35 = 1 \times 21 + 14$ \quad $\rightarrow$ \quad GCD(21, 14)

    $21 = 1 \times 14 + 7$ \quad $\rightarrow$ \quad GCD(14, 7) \qquad $14 = 2 \times 7 + 0$ \quad $\rightarrow$ \quad \highlight{GCD = 7}
  \end{exambox}
\end{frame}

% ===========================================================================
%  SECTION 4: สมภาค (Modular Arithmetic)
% ===========================================================================
\section{สมภาค (Modular Arithmetic)}

\begin{frame}{สมภาคคืออะไร?}
  \begin{explanation}[นิยาม]
    $a \equiv b \pmod{m}$ \quad ก็ต่อเมื่อ \quad $m \mid (a - b)$

    \medskip
    \textbf{อ่านว่า}: ``$a$ สมภาคกับ $b$ มอดุโล $m$'' (หรือ ``$a$ congruent to $b$ modulo $m$'')

    \textbf{ความหมาย}: $a$ และ $b$ \highlight{เหลือเศษเท่ากัน} เมื่อหารด้วย $m$
  \end{explanation}

  \begin{exambox}[ตัวอย่าง]
    $17 \equiv 5 \pmod{12}$ \quad เพราะ \quad $17 - 5 = 12$ หารด้วย $12$ ลงตัว

    $23 \equiv 3 \pmod{5}$ \quad เพราะ \quad $23 = 4\times5 + 3$ และ $3$ หารด้วย $5$ เหลือเศษ $3$

    $-7 \equiv 3 \pmod{5}$ \quad เพราะ \quad $-7 - 3 = -10$ หารด้วย $5$ ลงตัว
  \end{exambox}

  \begin{intuition}[คิดง่าย ๆ]
    ``mod $m$'' คือการวนรอบแบบนาฬิกา: นาฬิกามี $12$ ชั่วโมง
    — $14$ นาฬิกา $\equiv$ $2$ นาฬิกา (mod $12$)
  \end{intuition}
\end{frame}

\begin{frame}{สมบัติของสมภาค}
  \begin{explanation}[สมบัติพื้นฐาน]
    ถ้า $a \equiv b \pmod{m}$ และ $c \equiv d \pmod{m}$ แล้ว:
    \begin{enumerate}
      \item $a + c \equiv b + d \pmod{m}$ \quad (บวกกันได้)
      \item $a - c \equiv b - d \pmod{m}$ \quad (ลบกันได้)
      \item $ac \equiv bd \pmod{m}$ \quad (คูณกันได้)
      \item $a^n \equiv b^n \pmod{m}$ \quad (ยกกำลังได้)
    \end{enumerate}
  \end{explanation}

  \begin{exambox}[ตัวอย่าง: ลดทอนก่อนคำนวณ]
    จงหา $23 \times 18 \pmod{5}$

    \textbf{วิธีตรง:} $23 \times 18 = 414$, $414 \div 5 = 82$ เศษ $4$

    \textbf{วิธีสมภาค:} $23 \equiv 3 \pmod{5}$, $18 \equiv 3 \pmod{5}$ \qquad $3 \times 3 = 9 \equiv 4 \pmod{5}$ \quad \highlight{เท่ากัน!}
  \end{exambox}
\end{frame}

\begin{frame}{การหาเศษจากเลขยกกำลัง: ดูแพทเทิร์น}
  \begin{explanation}[แนวคิด]
    หา \textbf{วัฏจักร (cycle)} ของเศษจากการยกกำลัง:
    คำนวณ $a^1, a^2, a^3, \ldots \pmod{m}$ จนกว่าเศษจะซํ้า
  \end{explanation}

  \begin{exambox}[ตัวอย่าง: เศษของ $7^{2025}$ หารด้วย $10$]
    หลักหน่วยของ $7^n$ (mod $10$):

    $7^1 = 7 \equiv 7$ \quad
    $7^2 = 49 \equiv 9$ \quad
    $7^3 = 343 \equiv 3$ \quad
    $7^4 = 2401 \equiv 1$

    $7^5 \equiv 7$ \quad $\rightarrow$ \highlight{วัฏจักรยาว 4: $\{7, 9, 3, 1\}$}

    $2025 \div 4 = 506$ เศษ $1$ \quad $\rightarrow$ \quad $7^{2025} \equiv 7^1 \equiv 7 \pmod{10}$

    \highlight{เศษ $= 7$ (หลักหน่วยของ $7^{2025}$ คือ $7$)}
  \end{exambox}
\end{frame}

\begin{frame}[t]{ตัวอย่าง: เศษของเลขยกกำลังขนาดใหญ่ (ต่อ)}
  \begin{exambox}[โจทย์]
    $3^{100} \pmod{7}$
  \end{exambox}

  \vspace{\fill}

  \begin{exambox}[โจทย์]
    $2^{50} \pmod{13}$
  \end{exambox}

  \vspace{\fill}
\end{frame}

\begin{frame}[t]{ตัวอย่าง: ผลรวมเลขยกกำลังขนาดใหญ่}
  \begin{exambox}[โจทย์]
    เศษของ $2018^{2019} + 2019^{2020} + 2020^{2021}$ หารด้วย $13$
  \end{exambox}

  \vspace{\fill}
\end{frame}

\begin{frame}{ทฤษฎีบทของแฟร์มาต์ (Fermat's Little Theorem)}
  \begin{explanation}[Fermat's Little Theorem]
    ถ้า $p$ เป็นจำนวนเฉพาะ และ $p \nmid a$ แล้ว:

    \highlight{$a^{p-1} \equiv 1 \pmod{p}$}

    หรือรูปทั่วไป: $a^p \equiv a \pmod{p}$
  \end{explanation}

  \begin{exambox}[ตัวอย่าง]
    \textbf{1.} $2^{10} \equiv 1 \pmod{11}$ \quad (ตรวจสอบ: $1024 \div 11 = 93$ เศษ $1$)

    \textbf{2.} $5^{6} \equiv 1 \pmod{7}$ \quad (ตรวจสอบ: $15625 \div 7 = 2232$ เศษ $1$)

    \textbf{3.} หา $7^{100} \pmod{5}$:

    $7^4 \equiv 1 \pmod{5}$ \quad (Fermat: $p=5$)

    $100 = 4 \times 25$ \quad $\rightarrow$ \quad
    $7^{100} = (7^4)^{25} \equiv 1^{25} \equiv 1 \pmod{5}$
  \end{exambox}
\end{frame}

\begin{frame}[t]{ตัวอย่าง: โจทย์วันในสัปดาห์}
  \begin{exambox}[โจทย์]
    วันนี้เป็นวันเสาร์ อยากทราบว่าอีก $1000$ วันข้างหน้าจะเป็นวันอะไร?
  \end{exambox}

  \vspace{\fill}
\end{frame}

\begin{frame}[t]{ตัวอย่างเพิ่มเติม: สมภาค}
  \begin{exambox}[โจทย์]
    จำนวนที่หารด้วย 5 เหลือเศษ 3 และหารด้วย 7 เหลือเศษ 2
  \end{exambox}

  \vspace{\fill}
\end{frame}

\begin{frame}[t]{ตัวอย่าง: หลักหน่วยของเลขยกกำลังซ้อน}
  \begin{exambox}[โจทย์]
    จงหาหลักหน่วยของ $7^{7^{7}}$
  \end{exambox}

  \vspace{\fill}

  \begin{exambox}[โจทย์]
    จงหาหลักหน่วยของ $2^{2024} + 3^{2024}$
  \end{exambox}

  \vspace{\fill}
\end{frame}

% ===========================================================================
%  SECTION 5: เลขฐาน
% ===========================================================================
\section{เลขฐาน (Number Bases)}

\begin{frame}{เลขฐาน}
  \begin{explanation}[แนวคิด]
    ตัวเลขที่เราใช้ในชีวิตประจำวันเป็น \textbf{เลขฐานสิบ} (decimal)
    แต่เราสามารถเขียนจำนวนในฐานอื่น ๆ ได้
  \end{explanation}

  \begin{columns}[T]
    \column{0.5\textwidth}
      \begin{exambox}[การเปลี่ยนฐาน $n$ เป็นฐานสิบ]
        $(a_k a_{k-1} \ldots a_1 a_0)_n$

        $= a_k n^k + a_{k-1} n^{k-1} + \cdots + a_1 n + a_0$

        \textbf{ตัวอย่าง:}

        $(1011)_2 = 1\cdot8 + 0\cdot4 + 1\cdot2 + 1 = 11$

        $(2A)_{16} = 2 \times 16 + 10 = 42$
      \end{exambox}
    \column{0.5\textwidth}
      \begin{exambox}[การเปลี่ยนฐานสิบเป็นฐาน $n$]
        หารด้วย $n$ ซํ้า ๆ เอาเศษมาเรียงจากล่างขึ้นบน

        \textbf{ตัวอย่าง:} $12$ เป็นฐาน $2$ \\
        $12 \div 2 = 6$ เศษ $0$

        $6 \div 2 = 3$ เศษ $0$

        $3 \div 2 = 1$ เศษ $1$

        $1 \div 2 = 0$ เศษ $1$

        $(12)_{10} = (1100)_2$
      \end{exambox}
  \end{columns}
\end{frame}

% ===========================================================================
%  SUMMARY
% ===========================================================================
\section{สรุป}

\begin{frame}{สรุปเนื้อหา}
  \begin{explanation}[สิ่งที่ได้เรียนรู้]
    \begin{columns}[T]
      \column{0.5\textwidth}
        \begin{itemize}
          \item \textbf{การหารลงตัว}: $a \mid b$, กฎการหาร
          \item \textbf{จำนวนเฉพาะ}: มีแค่ $1$ กับตัวมันเองที่หารลงตัว
          \item \textbf{จำนวนประกอบ}: เขียนเป็นผลคูณของจำนวนเฉพาะได้
          \item \textbf{GCD/LCM}: แยกตัวประกอบ, Euclidean algorithm
          \item \textbf{GCD $\times$ LCM = $ab$}
        \end{itemize}
      \column{0.5\textwidth}
        \begin{itemize}
          \item \textbf{สมภาค}: $a \equiv b \pmod{m} \iff m \mid (a-b)$
          \item \textbf{สมบัติ}: บวก ลบ คูณ ยกกำลัง ใน mod ได้
          \item \textbf{Fermat}: $a^{p-1} \equiv 1 \pmod{p}$ ($p$ เฉพาะ)
          \item \textbf{เลขยกกำลัง}: หาวัฏจักรของเศษ
          \item \textbf{เลขฐาน}: เปลี่ยนไป-กลับระหว่างฐานต่าง ๆ
        \end{itemize}
    \end{columns}
  \end{explanation}
  \begin{center}
    \large เอกสารนี้เป็นส่วนหนึ่งของเนื้อหาเตรียมสอบ \textbf{สอวน. คอมพิวเตอร์}\\
    Special thanks to \textbf{Pana W.} for the original English-version slide.
  \end{center}
\end{frame}

\end{document}
